%==============================================================================
%  GENERAL CONCLUSION  (~2 pages, unnumbered chapter)
%==============================================================================
\chapter*{General Conclusion and Perspectives}
\addcontentsline{toc}{chapter}{General Conclusion and Perspectives}
\markboth{General Conclusion}{General Conclusion}

This end-of-studies project set out to bridge the gap between training a machine
learning model and running it reliably in production, by designing and building a
complete, self-hosted \gls{mlops} platform. Carried out within \textbf{INSOMEA}
and organized as four Scrum sprints, the work delivered a system that automates the
entire machine-learning lifecycle from a single web interface.

\section*{Objectives Achieved}
The platform meets the objectives defined at the outset:
\begin{itemize}[leftmargin=1.4em]
  \item secure, token-based authentication and project management
        (Sprint~1);
  \item upload of arbitrary code and data, isolated execution on a Kubernetes
        cluster, and automatic capture of metrics, parameters, and artifacts
        (Sprint~2);
  \item a model registry with lifecycle stages, one-click deployment as a live
        prediction service, and an in-interface tester (Sprint~3);
  \item a live observability dashboard, a public key-authenticated prediction
        \gls{api}, diagnostics that make failures understandable, and an
        administration area for user and platform management (Sprint~4).
\end{itemize}
The whole system is containerized, runs identically on a local cluster and on
\gls{aks}, and is continuously deployed to a public URL through an automated
\gls{cicd} pipeline.

\section*{Skills Gained}
The project was an opportunity to develop both technical and methodological
skills: container orchestration with Kubernetes; the \gls{mlops} toolchain
(pipeline orchestration, experiment tracking, model serving); full-stack
development with an asynchronous Python backend and a modern Angular frontend;
cloud deployment and networking; and the disciplined application of the Scrum
framework to deliver software incrementally.

\section*{Limitations}
The current system is deliberately scoped as a single-tenant platform. It does not
yet enforce per-key rate limiting on the public endpoint, has no automated drift or
data-quality monitoring, and schedules only \gls{cpu} workloads. These are
reasonable boundaries for an end-of-studies project, but they mark where a
production hardening effort would begin.

\section*{Perspectives}
Several directions would extend the platform:
\begin{itemize}[leftmargin=1.4em]
  \item \textbf{Multi-tenancy} with organizations, teams, and role-based access
        control;
  \item \textbf{\gls{gpu} scheduling} for deep-learning workloads;
  \item \textbf{model monitoring}, prediction logging, data and concept drift
        detection, and alerting;
  \item \textbf{autoscaling} of model servers based on traffic, and
        \textbf{A/B testing} or canary releases between model versions;
  \item \textbf{rate limiting and quotas} on the public \gls{api}.
\end{itemize}

In conclusion, the project resulted in a functional, end-to-end \gls{mlops}
platform that demonstrably automates the machine-learning lifecycle, and provided a
deep, hands-on understanding of the technologies and practices that make machine
learning operational.
